\section{Separate transformation and authentication}
XOR makes reuse visible without any knowledge of the quarter-round. If two messages
use the same keystream bytes, then
\[
  (P_1\oplus K)\oplus(P_2\oplus K)=P_1\oplus P_2.
\]
The repeated \(K\) cancels. A caller therefore needs a nonce-allocation policy for
each key as well as correct counter arithmetic inside one message. Incrementing
counters correctly does not make repeating the same initial state for another
message safe.

The transformation also has no way to recognize modification. Flipping a bit in a
ciphertext byte flips that same bit in the recovered plaintext. A successful return
from \texttt{chacha\_xor} reports only that the local count and counter conditions
were valid and the transformation completed. It does not authenticate the sender or
the message. ChaCha20-Poly1305 is a separate authenticated construction; implementing
this chapter's stream cipher does not implement that construction.

The prototype uses a fixed sequence of word operations for each block, without
key-indexed lookup tables. This observation about its source is narrower than a
constant-time claim about compiled hardware execution. A simulator instruction count
does not measure caches, speculative execution, power or other physical leakage.
Use the exercise to study dataflow and storage. Production cryptographic software
requires the surrounding protocol and implementation review as well.
